\Chapter{12}

\Verse{1}
{
    \zA{话说凤姐正与平儿说话，只见有人回说：“瑞大爷来了。” 凤姐命：“请进来 罢。” 贾瑞见请，心中暗喜，见了凤姐，满面陪笑，连连问好。 凤姐儿也假意殷勤
        让坐让茶。 贾瑞见凤姐如此打扮，越发酥倒，因饧了眼问道：“二哥哥怎么还不回 来？” 凤姐道：“不知什么缘故。” 贾瑞笑道：“别是路上有人绊住了脚，舍不得
        回来了罢？” 凤姐道：“可知男人家见一个爱一个也是有的。” 贾瑞笑道：“嫂子 这话错了，我就不是这样人。”}
}
{
    \eA{Lady Feng, it must be noticed in continuation of our narrative, was just
        engaged in talking with P'ing Erh, when they heard some one announce that
        Mr. Jui had come. Lady Feng gave orders that he should be invited to step
        in, and Chia Jui perceiving that he had been asked to walk in was at heart
        elated at the prospect of seeing her. With a face beaming with smiles, Lady
        Feng inquired again and again how he was;}
}
{
}

\Verse{2}
{
    \zA{凤姐笑道：“像你这样的人能有几个呢，十个里也 挑不出一个来！” 贾瑞听了，喜的抓耳挠腮，又道：“嫂子天天也闷的很。” 凤姐
        道：“正是呢，只盼个人来说话解解闷儿。” 贾瑞笑道：“我倒天天闲着。 若天天 过来替嫂子解解闷儿，可好么？”
        凤姐笑道：“你哄我呢!你那里肯往我这里来？” 贾瑞道：“我在嫂子面前若有一句谎话，天打雷劈!只因素日闻得人说，嫂子是个
        利害人，在你跟前一点也错不得，所以唬住我了。}
}
{
    \eA{and, with simulated tenderness she further pressed him to take a seat and
        urged him to have a cup of tea. Chia Jui noticed how still more voluptuous
        lady Feng looked in her present costume, and, as his eyes burnt with love,
        "How is it," he inquired, "that my elder brother Secundus is not yet back?"
        "What the reason is I cannot tell," lady Feng said by way of reply.}
}
{
}

\Verse{3}
{
    \zA{我如今见嫂子是个有说有笑极疼 人的，我怎么不来？ 死了也情愿。” 凤姐笑道：“果然你是个明白人，比蓉儿兄弟 两个强远了。
        我看他那样清秀，只当他们心里明白，谁知竟是两个糊涂虫，一点不 知人心。” 贾瑞听这话，越发撞在心坎上，由不得又往前凑一凑，觑着眼看凤姐的荷包，
        又问：“戴着什么戒指？” 凤姐悄悄的道：“放尊重些，别叫丫头们看见了。” 贾 瑞如听纶音佛语一般，忙往后退。 凤姐笑道：“你该去了。”}
}
{
    \eA{"May it not be," Chia Jui smilingly insinuated, "that some fair damsel has
        got hold of him on the way, and that he cannot brook to tear himself from
        her to come home?" "That makes it plain that there are those among men who
        fall in love with any girl they cast their eyes on," hinted lady Feng. "Your
        remarks are, sister-in-law, incorrect, for I'm none of this kind!" Chia Jui
        explained smirkingly.}
}
{
}

\Verse{4}
{
    \zA{贾瑞道：“我再坐一 坐儿，好狠心的嫂子！” 凤姐儿又悄悄的道：“大天白日人来人往，你就在这里也 不方便。
        你且去，等到晚上起了更你来，悄悄的在西边穿堂儿等我。” 贾瑞听了， 如得珍宝，忙问道：“你别哄我。 但是那里人过的多，怎么好躲呢？” 凤姐道：“你
        只放心，我把上夜的小厮们都放了假，两边门一关，再没别人了。” 贾瑞听了，喜 之不尽，忙忙的告辞而去，心内以为得手。
        盼到晚上，果然黑地里摸入荣府，趁掩门时钻入穿堂。}
}
{
    \eA{"How many like you can there be!" rejoined lady Feng with a sarcastic smile;
        "in ten, not one even could be picked out!" When Chia Jui heard these words,
        he felt in such high glee that he rubbed his ears and smoothed his cheeks.
        "My sister-in-law," he continued, "you must of course be extremely lonely
        day after day."}
}
{
}

\Verse{5}
{
    \zA{果见漆黑无一人来往， 贾母那边去的门已倒锁了，只有向东的门未关。 贾瑞侧耳听着，半日不见人来。 忽 听咯噔一声，东边的门也关上了。
        贾瑞急的也不敢则声，只得悄悄出来，将门撼了 撼，关得铁桶一般。 此时要出去亦不能了，南北俱是大墙，要跳也无攀援。 这屋内
        又是过堂风，空落落的，现是腊月天气，夜又长，朔风凛凛，侵肌裂骨，一夜几乎 不曾冻死。 好容易盼到早晨，只见一个老婆子先将东门开了进来，去叫西门，贾瑞
        瞅他背着脸，一溜烟抱了肩跑出来。}
}
{
    \eA{"Indeed I am," observed lady Feng, "and I only wish some one would come and
        have a chat with me to break my dull monotony." "I daily have ample
        leisure," Chia Jui ventured with a simper, "and wouldn't it be well if I
        came every day to dispel your dulness, sister-in-law?" "You are simply
        fooling me," exclaimed lady Feng laughing. "It isn't likely you would wish
        to come over here to me?"}
}
{
}

\Verse{6}
{
    \zA{幸而天气尚早，人都未起，从后门一径跑回家 去。 原来贾瑞父母早亡，只有他祖父代儒教养。 那代儒素日教训最严，不许贾瑞多
        走一步，生怕他在外吃酒赌钱，有误学业。 今忽见他一夜不归，只料定他在外非饮 即赌，嫖娼宿妓，那里想到这段公案？ 因此也气了一夜。 贾瑞也捻着一把汗，少不
        得回来撒谎，只说：“往舅舅家去了，天黑了，留我住了一夜。” 代儒道：“自来 出门非禀我不敢擅出，如何昨日私自去了？ 据此也该打，何况是撒谎！”}
}
{
    \eA{"If in your presence, sister-in-law, I utter a single word of falsehood, may
        the thunder from heaven blast me!" protested Chia Jui. "It's only because I
        had all along heard people say that you were a dreadful person, and that you
        cannot condone even the slightest shortcoming committed in your presence,
        that I was induced to keep back by fear;}
}
{
}

\Verse{7}
{
    \zA{因此发狠， 按倒打了三四十板，还不许他吃饭，叫他跪在院内读文章，定要补出十天工课来方 罢。
        贾瑞先冻了一夜，又挨了打，又饿着肚子，跪在风地里念文章，其苦万状。 此时贾瑞邪心未改，再不想到凤姐捉弄他。 过了两日，得了空儿，仍找寻凤姐。
        凤姐故意抱怨他失信，贾瑞急的起誓。 凤姐因他自投罗网，少不的再寻别计令他知 改，故又约他道：“今日晚上，你别在那里了，你在我这房后小过道儿里头那间空
        屋子里等我。 ――可别冒撞了！” 贾瑞道：“果真么？”}
}
{
    \eA{but after seeing you, on this occasion, so chatty, so full of fun and most
        considerate to others, how can I not come? were it to be the cause of my
        death, I would be even willing to come!" "You're really a clever person,"
        lady Feng observed sarcastically. "And oh so much superior to both Chia Jung
        and his brother!}
}
{
}

\Verse{8}
{
    \zA{凤姐道：“你不信就别来！” 贾瑞道：“必来，必来!死也要来的。” 凤姐道：“这会子你先去罢。” 贾瑞料定 晚间必妥，此时先去了。
        凤姐在这里便点兵派将，设下圈套。 那贾瑞只盼不到晚，偏偏家里亲戚又来了，吃了晚饭才去，那天已有掌灯时候；
        又等他祖父安歇，方溜进荣府，往那夹道中屋子里来等着，热锅上蚂蚁一般。 只是 左等不见人影，右听也没声响，心中害怕，不住猜疑道：“别是不来了，又冻我一
        夜不成？” 正自胡猜，只见黑？？ 的进来一个人。}
}
{
    \eA{Handsome as their presence was to look at, I imagined their minds to be full
        of intelligence, but who would have thought that they would, after all, be a
        couple of stupid worms, without the least notion of human affection!" The
        words which Chia Jui heard, fell in so much the more with his own
        sentiments, that he could not restrain himself from again pressing forward
        nearer to her;}
}
{
}

\Verse{9}
{
    \zA{贾瑞便打定是凤姐，不管青红皂 白，那人刚到面前，便如饿虎扑食、猫儿捕鼠的一般抱住，叫道：“亲嫂子，等死 我了！”
        说着，抱到屋里炕上就亲嘴扯裤子，满口里“亲爹”“亲娘”的乱叫起来。 那人只不做声，贾瑞便扯下自己的裤子来，硬帮帮就想顶入。 忽然灯光一闪，只见
        贾蔷举着个蜡台，照道：“谁在这屋里呢？” 只见炕上那人笑道：“瑞大叔要？ 我 呢！” 贾瑞不看则已，看了时真臊的无地可入。 你道是谁？ 却是贾蓉。}
}
{
    \eA{and as with eyes strained to give intentness to his view, he gazed at lady
        Feng's purse: "What rings have you got on?" he went on to ask. "You should
        be a little more deferential," remonstrated lady Feng in a low tone of
        voice, "so as not to let the waiting-maids detect us." Chia Jui withdrew
        backward with as much alacrity as if he had received an Imperial decree or a
        mandate from Buddha. "You ought to be going!"}
}
{
}

\Verse{10}
{
    \zA{贾瑞回身要跑， 被贾蔷一把揪住道：“别走!如今琏二婶子已经告到太太跟前，说你调戏他，他暂 时稳住你在这里。 太太听见气死过去了，这会子叫我来拿你。
        快跟我走罢！” 贾瑞 听了，魂不附体，只说：“好侄儿!你只说没有我，我明日重重的谢你！” 贾蔷道： “放你不值什么，只不知你谢我多少？
        况且口说无凭，写一张文契才算。” 贾瑞道： “这怎么落纸呢？” 贾蔷道：“这也不妨，写个赌钱输了，借银若干两，就完了。” 贾瑞道：“这也容易。”}
}
{
    \eA{lady Feng suggested, as she gave him a smile. "Do let me stay a while
        longer," entreated Chia Jui, "you are indeed ruthless, my sister-in-law."
        But with gentle voice did lady Feng again expostulate. "In broad daylight,"
        she said, "with people coming and going, it is not really convenient that
        you should abide in here;}
}
{
}

\Verse{11}
{
    \zA{贾蔷翻身出来，纸笔现成，拿来叫贾瑞写。 他两个做好做 歹，只写了五十两银子，画了押，贾蔷收起来。 然后撕掳贾蓉。 贾蓉先咬定牙不依，
        只说：“明日告诉族中的人评评理。” 贾瑞急的至于磕头。 贾蔷做好做歹的，也写 了一张五十两欠契才罢。 贾蔷又道：“如今要放你，我就担着不是。
        老太太那边的 门早已关了。 老爷正在厅上看南京来的东西，那一条路定难过去。 如今只好走后门。 要这一走，倘或遇见了人，连我也不好。
        等我先去探探，再来领你。}
}
{
    \eA{so you had better go, and when it's dark and the watch is set, you can come
        over, and quietly wait for me in the corridor on the Eastern side!" At these
        words, Chia Jui felt as if he had received some jewel or precious thing.
        "Don't make fun of me!" he remarked with vehemence. "The only thing is that
        crowds of people are ever passing from there, and how will it be possible
        for me to evade detection?"}
}
{
}

\Verse{12}
{
    \zA{这屋里你还藏 不住，少时就来堆东西，等我寻个地方。” 说毕，拉着贾瑞，仍息了灯，出至院外， 摸着大台阶底下，说道：“这窝儿里好。 只蹲着，别哼一声。
        等我来再走。” 说毕， 二人去了。 贾瑞此时身不由己，只得蹲在那台阶下。 正要盘算，只听头顶上一声响，哗喇
        喇一净桶尿粪从上面直泼下来，可巧浇了他一身一头。 贾瑞掌不住“嗳哟”一声， 忙又掩住口，不敢声张，满头满脸皆是尿屎，浑身冰冷打战。 只见贾蔷跑来叫：“快
        走，快走！”}
}
{
    \eA{"Set your mind at ease!" lady Feng advised; "I shall dismiss on leave all
        the youths on duty at night; and when the doors, on both sides, are closed,
        there will be no one else to come in!" Chia Jui was delighted beyond measure
        by the assurance, and with impetuous haste, he took his leave and went off;
        convinced at heart of the gratification of his wishes.}
}
{
}

\Verse{13}
{
    \zA{贾瑞方得了命，三步两步从后门跑到家中，天已三更，只得叫开了门。 家人见他这般光景，问：“是怎么了？” 少不得撒谎说：“天黑了，失脚掉在茅厕 里了。”
        一面即到自己房中更衣洗濯。 心下方想到凤姐玩他，因此发一回狠。 再想 想凤姐的模样儿标致，又恨不得一时搂在怀里。 胡思乱想，一夜也不曾合眼。 自此
        虽想凤姐，只不敢往荣府去了。 贾蓉等两个常常来要银子，他又怕祖父知道。 正是相思尚且难禁，况又添了债 务，日间工课又紧；}
}
{
    \eA{He continued, up to the time of dusk, a prey to keen expectation; and, when
        indeed darkness fell, he felt his way into the Jung mansion, availing
        himself of the moment, when the doors were being closed, to slip into the
        corridor, where everything was actually pitch dark, and not a soul to be
        seen going backwards or forwards.}
}
{
}

\Verse{14}
{
    \zA{他二十来岁的人，尚未娶亲，想着凤姐不得到手，自不免有些 “指头儿告了消乏”； 更兼两回冻恼奔波：因此三五下里夹攻，不觉就得了一病：
        心内发膨胀，口内无滋味，脚下如绵，眼中似醋，黑夜作烧，白日常倦，下溺遗精， 嗽痰带血，诸如此症，不上一年都添全了。 于是不能支持，一头躺倒，合上眼还只
        梦魂颠倒，满口胡话，惊怖异常。 百般请医疗治，诸如肉桂、附子、鳖甲、麦冬、 玉竹等药吃了有几十斤下去，也不见个动静。 倏又腊尽春回，这病更加沉重。}
}
{
    \eA{The door leading over to dowager lady Chia's apartments had already been put
        under key, and there was but one gate, the one on the East, which had not as
        yet been locked. Chia Jui lent his ear, and listened for ever so long, but
        he saw no one appear. Suddenly, however, was heard a sound like "lo teng,"
        and the east gate was also bolted;}
}
{
}

\Verse{15}
{
    \zA{代儒也着了忙，各处请医疗治，皆不见效。 因 后来吃“独参汤”，代儒如何有这力量，只得往荣府里来寻。 王夫人命凤姐秤二两 给他。
        凤姐回说：“前儿新近替老太太配了药，那整的太太又说留着送杨提督的太 太配药，偏偏昨儿我已经叫人送了去了。” 王夫人道：“就是咱们这边没了，你叫
        个人往你婆婆那里问问，或是你珍大哥哥那里有，寻些来凑着给人家。 吃好了，救 人一命，也是你们的好处。” 凤姐应了，也不遣人去寻。}
}
{
    \eA{but though Chia Jui was in a great state of impatience, he none the less did
        not venture to utter a sound. All that necessity compelled him to do was to
        issue, with quiet steps, from his corner, and to try the gates by pushing;
        but they were closed as firmly as if they had been made fast with iron
        bolts;}
}
{
}

\Verse{16}
{
    \zA{只将些渣末凑了几钱，命 人送去，只说：“太太叫送来的，再也没了。” 然后向王夫人说：“都寻了来了， 共凑了二两多，送去了。”
        那贾瑞此时要命心急，无药不吃，只是白花钱不见效。 忽然这日有个跛足道人 来化斋，口称专治冤孽之症。 贾瑞偏偏在内听见了，直着声叫喊，说：“快去请进
        那位菩萨来救命！” 一面在枕头上磕头。 众人只得带进那道士来。 贾瑞一把拉住， 连叫“菩萨救我！” 那道士叹道：“你这病非药可医。}
}
{
    \eA{and much though he may, at this juncture, have wished to find his way out,
        escape was, in fact, out of the question; on the south and north was one
        continuous dead wall, which, even had he wished to scale, there was nothing
        which he could clutch and pull himself up by. This room, besides, was one
        the interior (of which was exposed) to the wind, which entered through (the
        fissure) of the door;}
}
{
}

\Verse{17}
{
    \zA{我有个宝贝与你，你天天看 时，此命可保矣。” 说毕，从搭裢中取出个正面反面皆可照人的镜子来，――背上
        錾着“风月宝鉴”四字，――递与贾瑞道：“这物出自太虚幻境空灵殿上，警幻仙 子所制，专治邪思妄动之症，有济世保生之功。 所以带他到世上来，单与那些聪明
        俊秀、风雅王孙等照看。 千万不可照正面，只照背面，要紧，要紧!三日后我来收 取，管叫你病好。” 说毕，徉长而去。 众人苦留不住。
        贾瑞接了镜子，想道：“这道士倒有意思，我何不照一照试试？”}
}
{
    \eA{and was perfectly empty and bare; and the weather being, at this time, that
        of December, and the night too very long, the northerly wind, with its
        biting gusts, was sufficient to penetrate the flesh and to cleave the bones,
        so that the whole night long he had a narrow escape from being frozen to
        death;}
}
{
}

\Verse{18}
{
    \zA{想毕，拿起 那“宝鉴”来，向反面一照。 只见一个骷髅儿，立在里面。 贾瑞忙掩了，骂那道士： “混帐!如何吓我!我倒再照照正面是什么？”
        想着，便将正面一照，只见凤姐站在 里面点手儿叫他。 贾瑞心中一喜，荡悠悠觉得进了镜子，与凤姐云雨一番，凤姐仍 送他出来。
        到了床上，“嗳哟”了一声，一睁眼，镜子从新又掉过来，仍是反面立 着一个骷髅。 贾瑞自觉汗津津的，底下已遗了一滩精。}
}
{
    \eA{and he was yearning, with intolerable anxiety for the break of day, when he
        espied an old matron go first and open the door on the East side, and then
        come in and knock at the western gate. Chia Jui seeing that she had turned
        her face away, bolted out, like a streak of smoke, as he hugged his
        shoulders with his hands (from intense cold.)}
}
{
}

\Verse{19}
{
    \zA{心中到底不足，又翻过正面 来，只见凤姐还招手叫他，他又进去：如此三四次。 到了这次，刚要出镜子来，只 见两个人走来，拿铁锁把他套住，拉了就走。
        贾瑞叫道：“让我拿了镜子再走――” 只说这句就再不能说话了。 旁边伏侍的人只见他先还拿着镜子照，落下来，仍睁开眼拾在手内，末后镜子 掉下来，便不动了。
        众人上来看时，已经咽了气了，身子底下冰凉精湿遗下了一大 滩精。 这才忙着穿衣抬床。 代儒夫妇哭的死去活来，大骂道士：“是何妖道！”}
}
{
    \eA{As luck would have it, the hour was as yet early, so that the inmates of the
        house had not all got out of bed; and making his escape from the postern
        door, he straightaway betook himself home, running back the whole way. Chia
        Jui's parents had, it must be explained, departed life at an early period,
        and he had no one else, besides his grandfather Tai-ju, to take charge of
        his support and education.}
}
{
}

\Verse{20}
{
    \zA{遂 命人架起火来烧那镜子。 只听空中叫道：“谁叫他自己照了正面呢!你们自己以假 为真，为何烧我此镜？” 忽见那镜从房中飞出。
        代儒出门看时，却还是那个跛足道 人，喊道：“还我的风月宝鉴来！” 说着，抢了镜子，眼看着他飘然去了。 当下代儒没法，只得料理丧事，各处去报。
        三日起经，七日发引，寄灵铁槛寺 后。 一时贾家众人齐来吊问。 荣府贾赦赠银二十两，贾政也是二十两，宁府贾珍亦
        有二十两，其馀族中人贫富不一，或一二两、三四两不等。}
}
{
    \eA{This Tai-ju had, all along, exercised a very strict control, and would not
        allow Chia Jui to even make one step too many, in the apprehension that he
        might gad about out of doors drinking and gambling, to the neglect of his
        studies.}
}
{
}

\Verse{21}
{
    \zA{外又有各同窗家中分资， 也凑了二三十两。 代儒家道虽然淡薄，得此帮助，倒也丰丰富富完了此事。 谁知这年冬底，林如海因为身染重疾，写书来特接黛玉回去。
        贾母听了，未免 又加忧闷，只得忙忙的打点黛玉起身。 宝玉大不自在，争奈父女之情，也不好拦阻。 于是贾母定要贾琏送他去，仍叫带回来。
        一应土仪盘费，不消絮说，自然要妥贴的。 作速择了日期，贾琏同着黛玉辞别了众人，带领仆从，登舟往扬州去了。 要知端的，且听下回分解。 (本章完)}
}
{
    \eA{Seeing, on this unexpected occasion, that he had not come home the whole
        night, he simply felt positive, in his own mind, that he was certain to have
        run about, if not drinking, at least gambling, and dissipating in houses of
        the demi-monde up to the small hours; but he never even gave so much as a
        thought to the possibility of a public scandal, as that in which he was
        involved.}
}
{
}

\Verse{22}
{
    \zA{}
}
{
    \eA{The consequence was that during the whole length of the night he boiled with
        wrath. Chia Jui himself, on the other hand, was (in such a state of
        trepidation) that he could wipe the perspiration (off his face) by handfuls;
        and he felt constrained on his return home, to have recourse to deceitful
        excuses, simply explaining that he had been at his eldest maternal uncle's
        house, and that when it got dark, they kept him to spend the night there.}
}
{
}

\Verse{23}
{
    \zA{}
}
{
    \eA{"Hitherto," remonstrated Tai-ju, "when about to go out of doors, you never
        ventured to go, on your own hook, without first telling me about it, and how
        is it that yesterday you surreptitiously left the house? for this offence
        alone you deserve a beating, and how much more for the lie imposed upon me."
        Into such a violent fit of anger did he consequently fly that laying hands
        on him, he pulled him over and administered to him thirty or forty blows
        with a cane. Nor would he allow him to have anything to eat, but bade him
        remain on his knees in the court conning essays; impressing on his mind that
        he would not let him off, before he had made up for the last ten days'
        lessons.}
}
{
}

\Verse{24}
{
    \zA{}
}
{
    \eA{Chia Jui had in the first instance, frozen the whole night, and, in the next
        place, came in for a flogging. With a stomach, besides, gnawed by the pangs
        of hunger, he had to kneel in a place exposed to drafts reading the while
        literary compositions, so that the hardships he had to endure were of
        manifold kinds. Chia Jui's infamous intentions had at this junction
        undergone no change; but far from his thoughts being even then any idea that
        lady Feng was humbugging him, he seized, after the lapse of a couple of
        days, the first leisure moments to come again in search of that lady. Lady
        Feng pretended to bear him a grudge for his breach of faith, and Chia Jui
        was so distressed that he tried by vows and oaths (to establish his
        innocence.)}
}
{
}

\Verse{25}
{
    \zA{}
}
{
    \eA{Lady Feng perceiving that he had, of his own accord, fallen into the meshes
        of the net laid for him, could not but devise another plot to give him a
        lesson and make him know what was right and mend his ways. With this
        purpose, she gave him another assignation. "Don't go over there," she said,
        "to-night, but wait for me in the empty rooms giving on to a small passage
        at the back of these apartments of mine. But whatever you do, mind don't be
        reckless." "Are you in real earnest?" Chia Jui inquired. "Why, who wants to
        play with you?" replied lady Feng; "if you don't believe what I say, well
        then don't come!" "I'll come, I'll come, yea I'll come, were I even to die!"
        protested Chia Jui. "You should first at this very moment get away!"}
}
{
}

\Verse{26}
{
    \zA{}
}
{
    \eA{lady Feng having suggested, Chia Jui, who felt sanguine that when evening
        came, success would for a certainty crown his visit, took at once his
        departure in anticipation (of his pleasure.) During this interval lady Feng
        hastily set to work to dispose of her resources, and to add to her
        stratagems, and she laid a trap for her victim; while Chia Jui, on the other
        hand, was until the shades of darkness fell, a prey to incessant
        expectation. As luck would have it a relative of his happened to likewise
        come on that very night to their house and to only leave after he had dinner
        with them, and at an hour of the day when the lamps had already been lit;}
}
{
}

\Verse{27}
{
    \zA{}
}
{
    \eA{but he had still to wait until his grandfather had retired to rest before he
        could, at length with precipitate step, betake himself into the Jung
        mansion. Straightway he came into the rooms in the narrow passage, and
        waited with as much trepidation as if he had been an ant in a hot pan. He
        however waited and waited, but he saw no one arrive; he listened but not
        even the sound of a voice reached his ear. His heart was full of intense
        fear, and he could not restrain giving way to surmises and suspicion. "May
        it not be," he thought, "that she is not coming again; and that I may have
        once more to freeze for another whole night?" While indulging in these
        erratic reflections, he discerned some one coming, looking like a black
        apparition, who Chia Jui readily concluded, in his mind, must be lady Feng;}
}
{
}

\Verse{28}
{
    \zA{}
}
{
    \eA{so that, unmindful of distinguishing black from white, he as soon as that
        person arrived in front of him, speedily clasped her in his embrace, like a
        ravenous tiger pouncing upon its prey, or a cat clawing a rat, and cried:
        "My darling sister, you have made me wait till I'm ready to die." As he
        uttered these words, he dragged the comer, in his arms, on to the couch in
        the room; and while indulging in kisses and protestations of warm love, he
        began to cry out at random epithets of endearment. Not a sound, however,
        came from the lips of the other person;}
}
{
}

\Verse{29}
{
    \zA{}
}
{
    \eA{and Chia Jui had in the fulness of his passion, exceeded the bounds of timid
        love and was in the act of becoming still more affectionate in his
        protestations, when a sudden flash of a light struck his eye, by the rays of
        which he espied Chia Se with a candle in hand, casting the light round the
        place, "Who's in this room?" he exclaimed. "Uncle Jui," he heard some one on
        the couch explain, laughing, "was trying to take liberties with me!" Chia
        Jui at one glance became aware that it was no other than Chia Jung; and a
        sense of shame at once so overpowered him that he could find nowhere to hide
        himself; nor did he know how best to extricate himself from the dilemma.}
}
{
}

\Verse{30}
{
    \zA{}
}
{
    \eA{Turning himself round, he made an attempt to make good his escape, when Chia
        Se with one grip clutched him in his hold. "Don't run away," he said;
        "sister-in-law Lien has already reported your conduct to madame Wang; and
        explained that you had tried to make her carry on an improper flirtation
        with you; that she had temporised by having recourse to a scheme to escape
        your importunities, and that she had imposed upon you in such a way as to
        make you wait for her in this place. Our lady was so terribly incensed, that
        she well-nigh succumbed; and hence it is that she bade me come and catch
        you! Be quick now and follow me, and let us go and see her."}
}
{
}

\Verse{31}
{
    \zA{}
}
{
    \eA{After Chia Jui had heard these words, his very soul could not be contained
        within his body. "My dear nephew," he entreated, "do tell her that it wasn't
        I; and I'll show you my gratitude to-morrow in a substantial manner."
        "Letting you off," rejoined Chia Se, "is no difficult thing; but how much, I
        wonder, are you likely to give? Besides, what you now utter with your lips,
        there will be no proof to establish; so you had better write a promissory
        note." "How could I put what happened in black and white on paper?" observed
        Chia Jui. "There's no difficulty about that either!" replied Chia Se; "just
        write an account of a debt due, for losses in gambling, to some one outside;
        for payment of which you had to raise funds, by a loan of a stated number of
        taels, from the head of the house; and that will be all that is required."}
}
{
}

\Verse{32}
{
    \zA{}
}
{
    \eA{"This is, in fact, easy enough!" Chia Jui having added by way of answer;
        Chia Se turned round and left the room; and returning with paper and
        pencils, which had been got ready beforehand for the purpose, he bade Chia
        Jui write. The two of them (Chia Jung and Chia Se) tried, the one to do a
        good turn, and the other to be perverse in his insistence; but (Chia Jui)
        put down no more than fifty taels, and appended his signature. Chia Se
        pocketed the note, and endeavoured subsequently to induce Chia Jung to come
        away; but Chia Jung was, at the outset, obdurate and unwilling to give in,
        and kept on repeating; "To-morrow, I'll tell the members of our clan to look
        into your nice conduct!"}
}
{
}

\Verse{33}
{
    \zA{}
}
{
    \eA{These words plunged Chia Jui in such a state of dismay, that he even went so
        far as to knock his head on the ground; but, as Chia Se was trying to get
        unfair advantage of him though he had at first done him a good turn, he had
        to write another promissory note for fifty taels, before the matter was
        dropped. Taking up again the thread of the conversation, Chia Se remarked,
        "Now when I let you go, I'm quite ready to bear the blame! But the gate at
        our old lady's over there is already bolted, and Mr. Chia Cheng is just now
        engaged in the Hall, looking at the things which have arrived from Nanking,
        so that it would certainly be difficult for you to pass through that way.
        The only safe course at present is by the back gate;}
}
{
}

\Verse{34}
{
    \zA{}
}
{
    \eA{but if you do go by there, and perchance meet any one, even I will be in for
        a mess; so you might as well wait until I go first and have a peep, when
        I'll come and fetch you! You couldn't anyhow conceal yourself in this room;
        for in a short time they'll be coming to stow the things away, and you had
        better let me find a safe place for you." These words ended, he took hold of
        Chia Jui, and, extinguishing again the lantern, he brought him out into the
        court, feeling his way up to the bottom of the steps of the large terrace.
        "It's safe enough in this nest," he observed, "but just squat down quietly
        and don't utter a sound; wait until I come back before you venture out."
        Having concluded this remark, the two of them (Chia Se and Chia Jung) walked
        away;}
}
{
}

\Verse{35}
{
    \zA{}
}
{
    \eA{while Chia Jui was, all this time, out of his senses, and felt constrained
        to remain squatting at the bottom of the terrace stairs. He was about to
        consider what course was open for him to adopt, when he heard a noise just
        over his head; and, with a splash, the contents of a bucket, consisting
        entirely of filthy water, was emptied straight down over him from above,
        drenching, as luck would have it, his whole person and head. Chia Jui could
        not suppress an exclamation. "Ai ya!" he cried, but he hastily stopped his
        mouth with his hands, and did not venture to give vent to another sound.}
}
{
}

\Verse{36}
{
    \zA{}
}
{
    \eA{His whole head and face were a mass of filth, and his body felt icy cold.
        But as he shivered and shook, he espied Chia Se come running. "Get off," he
        shouted, "with all speed! off with you at once!" As soon as Chia Jui
        returned to life again, he bolted with hasty strides, out of the back gate,
        and ran the whole way home. The night had already reached the third watch,
        so that he had to knock at the door for it to be opened. "What's the
        matter?" inquired the servants, when they saw him in this sorry plight; (an
        inquiry) which placed him in the necessity of making some false excuse. "The
        night was dark," he explained, "and my foot slipped and I fell into a
        gutter." Saying this, he betook himself speedily to his own apartment;}
}
{
}

\Verse{37}
{
    \zA{}
}
{
    \eA{and it was only after he had changed his clothes and performed his
        ablutions, that he began to realise that lady Feng had made a fool of him.
        He consequently gave way to a fit of wrath; but upon recalling to mind the
        charms of lady Feng's face, he felt again extremely aggrieved that he could
        not there and then clasp her in his embrace, and as he indulged in these
        wild thoughts and fanciful ideas, he could not the whole night long close
        his eyes. From this time forward his mind was, it is true, still with lady
        Feng, but he did not have the courage to put his foot into the Jung mansion;}
}
{
}

\Verse{38}
{
    \zA{}
}
{
    \eA{and with Chia Jung and Chia Se both coming time and again to dun him for the
        money, he was likewise full of fears lest his grandfather should come to
        know everything. His passion for lady Feng was, in fact, already a burden
        hard to bear, and when, moreover, the troubles of debts were superadded to
        his tasks, which were also during the whole day arduous, he, a young man of
        about twenty, as yet unmarried, and a prey to constant cravings for lady
        Feng, which were difficult to gratify, could not avoid giving way, to a
        great extent, to such evil habits as exhausted his energies.}
}
{
}

\Verse{39}
{
    \zA{}
}
{
    \eA{His lot had, what is more, been on two occasions to be frozen, angered and
        to endure much hardship, so that with the attacks received time and again
        from all sides, he unconsciously soon contracted an organic disease. In his
        heart inflammation set in; his mouth lost the sense of taste; his feet got
        as soft as cotton from weakness; his eyes stung, as if there were vinegar in
        them. At night, he burnt with fever. During the day, he was repeatedly under
        the effects of lassitude. Perspiration was profuse, while with his
        expectorations of phlegm, he brought up blood. The whole number of these
        several ailments came upon him, before the expiry of a year, (with the
        result that) in course of time, he had not the strength to bear himself up.}
}
{
}

\Verse{40}
{
    \zA{}
}
{
    \eA{Of a sudden, he would fall down, and with his eyes, albeit closed, his
        spirit would be still plunged in confused dreams, while his mouth would be
        full of nonsense and he would be subject to strange starts. Every kind of
        doctor was asked to come in, and every treatment had recourse to; and,
        though of such medicines as cinnamon, aconitum seeds, turtle shell,
        ophiopogon, Yü-chü herb, and the like, he took several tens of catties, he
        nevertheless experienced no change for the better; so that by the time the
        twelfth moon drew once again to an end, and spring returned, this illness
        had become still more serious.}
}
{
}

\Verse{41}
{
    \zA{}
}
{
    \eA{Tai-ju was very much concerned, and invited doctors from all parts to attend
        to him, but none of them could do him any good. And as later on, he had to
        take nothing else but decoctions of pure ginseng, Tai-ju could not of course
        afford it. Having no other help but to come over to the Jung mansion, and
        make requisition for some, Madame Wang asked lady Feng to weigh two taels of
        it and give it to him. "The other day," rejoined lady Feng, "not long ago,
        when we concocted some medicine for our dowager lady, you told us, madame,
        to keep the pieces that were whole, to present to the spouse of General Yang
        to make physic with, and as it happens it was only yesterday that I sent
        some one round with them."}
}
{
}

\Verse{42}
{
    \zA{}
}
{
    \eA{"If there's none over here in our place," suggested madame Wang, "just send
        a servant to your mother-in-law's, on the other side, to inquire whether
        they have any. Or it may possibly be that your elder brother-in-law Chen,
        over there, might have a little. If so, put all you get together, and give
        it to them; and when he shall have taken it, and got well and you shall have
        saved the life of a human being, it will really be to the benefit of you
        all." Lady Feng acquiesced; but without directing a single person to
        institute any search, she simply took some refuse twigs, and making up a few
        mace, she despatched them with the meagre message that they had been sent by
        madame Wang, and that there was, in fact, no more;}
}
{
}

\Verse{43}
{
    \zA{}
}
{
    \eA{subsequently reporting to madame Wang that she had asked for and obtained
        all there was and that she had collected as much as two taels, and forwarded
        it to them. Chia Jui was, meanwhile, very anxious to recover his health, so
        that there was no medicine that he would not take, but the outlay of money
        was of no avail, for he derived no benefit. On a certain day and at an
        unexpected moment, a lame Taoist priest came to beg for alms, and he averred
        that he had the special gift of healing diseases arising from grievances
        received, and as Chia Jui happened, from inside, to hear what he said, he
        forthwith shouted out: "Go at once, and bid that divine come in and save my
        life!" while he reverentially knocked his head on the pillow.}
}
{
}

\Verse{44}
{
    \zA{}
}
{
    \eA{The whole bevy of servants felt constrained to usher the Taoist in; and Chia
        Jui, taking hold of him with a dash, "My Buddha!" he repeatedly cried out,
        "save my life!" The Taoist heaved a sigh. "This ailment of yours," he
        remarked, "is not one that could be healed with any medicine; I have a
        precious thing here which I'll give you, and if you gaze at it every day,
        your life can be saved!" When he had done talking, he produced from his
        pouch a looking-glass which could reflect a person's face on the front and
        back as well. On the upper part of the back were engraved the four
        characters: "Precious Mirror of Voluptuousness."}
}
{
}

\Verse{45}
{
    \zA{}
}
{
    \eA{Handing it over to Chia Jui: "This object," he proceeded, "emanates from the
        primordial confines of the Great Void and has been wrought by the Monitory
        Dream Fairy in the Palace of Unreality and Spirituality, with the sole
        intent of healing the illnesses which originate from evil thoughts and
        improper designs. Possessing, as it does, the virtue of relieving mankind
        and preserving life, I have consequently brought it along with me into the
        world, but I only give it to those intelligent preëminent and refined
        princely men to set their eyes on. On no account must you look at the front
        side; and you should only gaze at the back of it; this is urgent, this is
        expedient! After three days, I shall come and fetch it away;}
}
{
}

\Verse{46}
{
    \zA{}
}
{
    \eA{by which time, I'm sure, it will have made him all right." These words
        finished, he walked away with leisurely step, and though all tried to detain
        him, they could not succeed. Chia Jui received the mirror. "This Taoist," he
        thought, "would seem to speak sensibly, and why should I not look at it and
        try its effect?" At the conclusion of these thoughts, he took up the Mirror
        of Voluptuousness, and cast his eyes on the obverse side; but upon
        perceiving nought else than a skeleton standing in it, Chia Jui sustained
        such a fright that he lost no time in covering it with his hands and in
        abusing the Taoist. "You good-for-nothing!" he exclaimed, "why should you
        frighten me so?}
}
{
}

\Verse{47}
{
    \zA{}
}
{
    \eA{but I'll go further and look at the front and see what it's like." While he
        reflected in this manner, he readily looked into the face of the mirror,
        wherein he caught sight of lady Feng standing, nodding her head and
        beckoning to him. With one gush of joy, Chia Jui felt himself, in a vague
        and mysterious manner, transported into the mirror, where he held an
        affectionate tête-à-tête with lady Feng. Lady Feng escorted him out again.
        On his return to bed, he gave vent to an exclamation of "Ai yah!" and
        opening his eyes, he turned the glass over once more; but still, as
        hitherto, stood the skeleton in the back part. Chia Jui had, it is true,
        experienced all the pleasant sensations of a tête-à-tête, but his heart
        nevertheless did not feel gratified;}
}
{
}

\Verse{48}
{
    \zA{}
}
{
    \eA{so that he again turned the front round, and gazed at lady Feng, as she
        still waved her hand and beckoned to him to go. Once more entering the
        mirror, he went on in the same way for three or four times, until this
        occasion, when just as he was about to issue from the mirror, he espied two
        persons come up to him, who made him fast with chains round the neck, and
        hauled him away. Chia Jui shouted. "Let me take the mirror and I'll come
        along." But only this remark could he utter, for it was forthwith beyond his
        power to say one word more.}
}
{
}

\Verse{49}
{
    \zA{}
}
{
    \eA{The servants, who stood by in attendance, saw him at first still holding the
        glass in his hand and looking in, and then, when it fell from his grasp,
        open his eyes again to pick it up, but when at length the mirror dropped,
        and he at once ceased to move, they in a body came forward to ascertain what
        had happened to him. He had already breathed his last. The lower part of his
        body was icy-cold; his clothes moist from profuse perspiration. With all
        promptitude they changed him there and then, and carried him to another bed.
        Tai-ju and his wife wept bitterly for him, to the utter disregard of their
        own lives, while in violent terms they abused the Taoist priest. "What kind
        of magical mirror is it?" they asked. "If we don't destroy this glass, it
        will do harm to not a few men in the world!"}
}
{
}

\Verse{50}
{
    \zA{}
}
{
    \eA{Having forthwith given directions to bring fire and burn it, a voice was
        heard in the air to say, "Who told you to look into the face of it? You
        yourselves have mistaken what is false for what is true, and why burn this
        glass of mine?" Suddenly the mirror was seen to fly away into the air; and
        when Tai-ju went out of doors to see, he found no one else than the limping
        Taoist, shouting, "Who is he who wishes to destroy the Mirror of
        Voluptuousness?" While uttering these words, he snatched the glass, and, as
        all eyes were fixed upon him, he moved away lissomely, as if swayed by the
        wind.}
}
{
}

\Verse{51}
{
    \zA{}
}
{
    \eA{Tai-ju at once made preparations for the funeral and went everywhere to give
        notice that on the third day the obsequies would commence, that on the
        seventh the procession would start to escort the coffin to the Iron Fence
        Temple, and that on the subsequent day, it would be taken to his original
        home. Not much time elapsed before all the members of the Chia family came,
        in a body, to express their condolences. Chia She, of the Jung Mansion,
        presented twenty taels, and Chia Cheng also gave twenty taels. Of the Ning
        Mansion, Chia Chen likewise contributed twenty taels. The remainder of the
        members of the clan, of whom some were poor and some rich, and not equally
        well off, gave either one or two taels, or three or four, some more, some
        less.}
}
{
}

\Verse{52}
{
    \zA{}
}
{
    \eA{Among strangers, there were also contributions, respectively presented by
        the families of his fellow-scholars, amounting, likewise, collectively to
        twenty or thirty taels. The private means of Tai-ju were, it is true,
        precarious, but with the monetary assistance he obtained, he anyhow
        performed the funeral rites with all splendour and éclat. But who would have
        thought it, at the close of winter of this year, Lin Ju-hai contracted a
        serious illness, and forwarded a letter, by some one, with the express
        purpose of fetching Lin Tai-yü back. These tidings, when they reached
        dowager lady Chia, naturally added to the grief and distress (she already
        suffered), but she felt compelled to make speedy preparations for Tai-yü's
        departure.}
}
{
}

\Verse{53}
{
    \zA{}
}
{
    \eA{Pao-yü too was intensely cut up, but he had no alternative but to defer to
        the affection of father and daughter; nor could he very well place any
        hindrance in the way. Old lady Chia, in due course, made up her mind that
        she would like Chia Lien to accompany her, and she also asked him to bring
        her back again along with him. But no minute particulars need be given of
        the manifold local presents and of the preparations, which were, of course,
        everything that could be wished for in excellence and perfectness.}
}
{
}

\Verse{54}
{
    \zA{}
}
{
    \eA{Forthwith the day for starting was selected, and Chia Lien, along with Lin
        Tai-yü, said good-bye to all the members of the family, and, followed by
        their attendants, they went on board their boats, and set out on their
        journey for Yang Chou. But, Reader, should you have any wish to know fuller
        details, listen to the account given in the subsequent Chapter.}
}
{
}
